\section{Концептуальна модель системи}

Обʼєкт моделювання -- інформаційна система, що виконує завдання типів A, B і~C на одній
обчислювальній системі зі спільною оперативною памʼяттю; завдання, які не можуть стартувати одразу,
чекають у черзі, порядок вибору з якої задає одна з пʼяти стратегій. Часові величини вимірюються у хвилинах (умовна одиниця модельного часу -- 1~хв),
інтенсивності -- у завданнях за хвилину.

\subsection{Склад системи та звʼязки між елементами}

Структура системи (рисунок~\ref{fig:structure}) містить такі елементи:
\begin{enumerate}
  \item три незалежні джерела завдань типів A, B і~C з інтенсивностями $\lambda_A$, $\lambda_B$, $\lambda_C$;
  \item черга, упорядкована за пріоритетами типів: три підчерги (по одній на тип) з FIFO
        всередині, довжина не обмежена, завдання не втрачаються;
  \item диспетчер, що обирає наступне завдання для запуску за чинною стратегією~$s$;
  \item обчислювальна підсистема з оперативною памʼяттю -- єдиним спільним обмеженим ресурсом.
\end{enumerate}

Умова задає лише середні інтенсивності, тож прийнято припущення про пуассонівські потоки з
інтервалами $\exp(1/\lambda_i)$ -- від багатьох незалежних користувачів; час виконання за умовою
експоненційний з~$\mu_i$.

Обмеження на одночасне виконання умова формулює так: завдання типу~A «можуть виконуватися
паралельно із завданням свого класу або із завданням типу~B», аналогічно для~B, а завдання
типу~C монополізують систему. Граничну кількість одночасних завдань прямо не вказано, тому
прийнято припущення: одночасно виконуються не більше двох завдань (пари AA, AB, BB), бо
в умові вжито однину -- «із завданням», тобто напарник щонайбільше один. Його відтворено обсягом
памʼяті $M=2$~умовні одиниці: A і~B займають по $v_A=v_B=1$~од., C -- $v_C=2$~од., тобто всю
памʼять (монополізація). Допустимі стани: порожня, A, B, AA, AB, BB, C. Наслідки іншого
тлумачення оцінює зміна~$M$ (розділи~3 і~5).

\begin{figure}
  \centering
  % Рисунок 1.1 — структура інформаційної системи (концептуальна модель)
\begin{tikzpicture}[
    font=\setstretch{1}\fontsize{11}{13}\selectfont, >={Stealth[length=2.2mm]}, node distance=4mm and 6mm,
    src/.style={draw, ellipse, align=center, inner sep=1.5pt, minimum width=22mm},
    queue/.style={draw, rectangle, minimum width=24mm, minimum height=9mm, align=left, text width=13mm},
    blk/.style={draw, rectangle, rounded corners=2pt, align=center, inner sep=4pt},
    flow/.style={->, thick},
    res/.style={->, dashed},
  ]
  % Джерела
  \node[src] (sA) {Джерело A\\$\lambda_A$};
  \node[src, below=of sA] (sB) {Джерело B\\$\lambda_B$};
  \node[src, below=of sB] (sC) {Джерело C\\$\lambda_C$};
  % Черги
  \node[queue, right=of sA] (qA) {Черга A};
  \node[queue, right=of sB] (qB) {Черга B};
  \node[queue, right=of sC] (qC) {Черга C};
  \foreach \q in {qA,qB,qC}{
    \foreach \x in {0.72,0.84,0.96}{
      \draw ($(\q.north west)!\x!(\q.north east)$) -- ($(\q.south west)!\x!(\q.south east)$);
    }
  }
  \node[draw, dotted, inner sep=3pt, fit=(qA)(qC), label={[align=center]above:Упорядкована\\черга}] (Q) {};
  % Диспетчер
  \node[blk, right=6mm of Q, minimum height=36mm, text width=24mm] (disp)
    {Вибір завдання за~стратегією пріоритетів~$s$};
  % Обробка
  \node[blk, right=7mm of disp, text width=36mm, minimum height=24mm] (cpu)
    {Виконання завдань\\[2pt]
     A: $v_A=1$, $1/\mu_A$\\
     B: $v_B=1$, $1/\mu_B$\\
     C: $v_C=2$, $1/\mu_C$};
  \node[blk, above=7mm of cpu, text width=36mm] (mem) {Оперативна памʼять\\$M=2$ од.};
  \node[src, minimum width=14mm, right=5mm of cpu] (out) {Вихід};
  % Зв'язки
  \draw[flow] (sA) -- (qA);
  \draw[flow] (sB) -- (qB);
  \draw[flow] (sC) -- (qC);
  \draw[flow] (qA.east) -- (qA.east -| disp.west);
  \draw[flow] (qB.east) -- (qB.east -| disp.west);
  \draw[flow] (qC.east) -- (qC.east -| disp.west);
  \draw[flow] (disp.east |- cpu) -- (cpu.west);
  \draw[flow] (cpu) -- (out);
  \draw[res] (mem.west) -| node[pos=0.25, above, font=\fontsize{10}{12}\selectfont] {захоплення $v_i$} (disp.north);
  \draw[res] ([xshift=6mm]cpu.north) -- node[right, font=\fontsize{10}{12}\selectfont] {звільнення $v_i$} ([xshift=6mm]mem.south);
\end{tikzpicture}

  \caption{Структура інформаційної системи обробки завдань}
  \label{fig:structure}
\end{figure}

Стратегії упорядкування з умови наведено в таблиці~\ref{tab:strategies}.

\begin{kpitable}{|L{3.4cm}|c|c|c|c|c|}{Стратегії упорядкування черги\label{tab:strategies}}
  {\textbf{Стратегія $s$} & \textbf{1} & \textbf{2} & \textbf{3} & \textbf{4} & \textbf{5}}
  Упорядкування & A $>$ B $>$ C & B $>$ A $>$ C & C $>$ B $>$ A & C $>$ A $>$ B & C $>$ (A $=$ B) \\ \hline
\end{kpitable}

\subsection{Опис процесу функціонування системи}\label{sec:process}

Завдання типу~$i$ стає в кінець своєї підчерги, починає виконуватися, щойно справджуються умови
старту (рисунок~\ref{fig:process}), а після виконання звільняє памʼять і виходить із системи.
Умови перевіряють після кожного надходження й звільнення памʼяті; за двох вільних одиниць можуть
одночасно стартувати два завдання A/B.

Умова визначає упорядкування саме в черзі, тож черга -- єдина впорядкована послідовність, яку
обслуговують з початку. Звідси правила:
\begin{itemize}
  \item резервування для завдань вищого пріоритету: завдання нижчого пріоритету не обходить
        завдання вищого пріоритету, що очікує, навіть за вільної памʼяті (інакше в стратегіях
        3--5 завдання~C майже не отримувало б усю памʼять); у стратегіях 1--2 C стартує лише за
        порожніх підчерг A і~B;
  \item обслуговування без переривання: пріоритет діє лише під час вибору наступного завдання;
  \item рівні пріоритети (стратегія~5): якщо на одну вільну одиницю памʼяті претендують
        завдання~A і~B, запускається одне з них з імовірностями $0{,}5$.
\end{itemize}

\begin{figure}
  \centering
  % Рисунок 1.2 — процес проходження одного завдання типу i через систему
\begin{tikzpicture}[
    font=\setstretch{1}\fontsize{11}{13}\selectfont, >={Stealth[length=2.2mm]}, node distance=3mm,
    ev/.style={draw, circle, minimum size=6mm},
    act/.style={draw, rectangle, rounded corners=3pt, align=center, text width=80mm, minimum height=6mm},
    dec/.style={draw, diamond, aspect=2.2, align=center, text width=30mm, inner sep=1pt},
    note/.style={draw, dashed, rectangle, align=left, text width=46mm, font=\setstretch{1}\fontsize{10}{12}\selectfont},
    flow/.style={->, thick},
  ]
  \node[ev] (start) {};
  \node[act, below=of start] (arr) {Надходження завдання типу $i\in\{A,B,C\}$};
  \node[act, below=of arr] (enq) {Постановка в кінець черги типу $i$};
  \node[dec, below=6mm of enq] (cond) {Умови старту виконано?};
  \node[act, below=6mm of cond] (take) {Захоплення $v_i$ од. памʼяті, вихід із черги};
  \node[act, below=of take] (serv) {Виконання протягом випадкового\\часу з~середнім $1/\mu_i$};
  \node[act, below=of serv] (rel) {Звільнення $v_i$ од. памʼяті};
  \node[act, below=of rel] (dep) {Вихід завдання із системи};
  \node[ev, very thick, below=of dep] (end) {};
  \node[act, text width=33mm, left=9mm of cond] (wait) {Очікування в черзі};
  \node[note, right=7mm of cond] (cn) {Умови старту:\\
    1) завдання перше в черзі типу $i$;\\
    2) у черзі немає завдань вищого пріоритету;\\
    3) вільна памʼять не менша за $v_i$};

  \draw[flow] (start) -- (arr);
  \draw[flow] (arr) -- (enq);
  \draw[flow] (enq) -- (cond);
  \draw[flow] (cond) -- node[right] {так} (take);
  \draw[flow] (cond.west) -- node[above] {ні} (wait.east);
  \draw[flow] (wait.north) |- ($(enq.south)!0.5!(cond.north)$);
  \draw[flow] (take) -- (serv);
  \draw[flow] (serv) -- (rel);
  \draw[flow] (rel) -- (dep);
  \draw[flow] (dep) -- (end);
  \draw[dotted] (cond.east) -- (cn.west);
  \draw[->, dashed] (rel.west) -| node[pos=0.75, left, align=right, font=\setstretch{1}\fontsize{10}{12}\selectfont]
    {звільнення\\памʼяті:\\повторна\\перевірка} (wait.south);
\end{tikzpicture}

  \caption{Процес проходження завдання типу~$i$ через систему}
  \label{fig:process}
\end{figure}

Спрощення: час виконання не залежить від паралельного виконання іншого завдання, затримки
диспетчеризації нехтовні.

За інтенсивностей з умови задачі система перевантажена: частка памʼяті, потрібна для
обслуговування всього потоку, становить
\begin{equation}\label{eq:rho}
  \rho = \frac{1}{M}\sum_{i\in\{A,B,C\}} \frac{\lambda_i v_i}{\mu_i}
       = \frac{1}{2}\left(\frac{0{,}2}{0{,}25} + \frac{0{,}07}{0{,}2}
          + \frac{0{,}06\cdot 2}{0{,}1}\right) = 1{,}175 > 1,
\end{equation}
тож черги зростають необмежено за будь-якої стратегії. Тому інтенсивності всіх типів
варіюються спільним коефіцієнтом~$\gamma$: $\lambda_i = \gamma\lambda_i^0$, де $\lambda_i^0$ --
значення з умови. Необхідна умова стійкості $\gamma\cdot 1{,}175<1$ дає $\gamma<0{,}85$; через
простоювання памʼяті під час резервування фактична межа нижча й залежить від стратегії (розділ~4).

\subsection{Вхідні змінні та параметри моделі}

Вхідні змінні наведено в таблиці~\ref{tab:inputs}, параметри -- у таблиці~\ref{tab:params};
$\exp(a)$ -- експоненційна випадкова величина із середнім~$a$~\cite{stetsenko-kursova}.

\begin{kpitable}{|L{4.6cm}|c|L{2.4cm}|L{2.7cm}|L{3.1cm}|}{Вхідні змінні моделі\label{tab:inputs}}
  {\textbf{Змінна} & \textbf{Позначення} & \textbf{Початкове значення} & \textbf{Діапазон} & \textbf{Закон розподілу}}
  Стратегія упорядкування черги & $s$ & 1 & $\{1,\dots,5\}$, табл.~\ref{tab:strategies} & детермінована \\ \hline
  Коефіцієнт інтенсивності надходження & $\gamma$ & 1 & $[0{,}4;\ 1{,}0]$ & детермінована \\ \hline
  Інтенсивності надходження A, B, C, завд./хв & $\lambda_A, \lambda_B, \lambda_C$ &
    $0{,}2$\newline $0{,}07$\newline $0{,}06$ &
    $[0{,}08;\ 0{,}2]$\newline $[0{,}028;\ 0{,}07]$\newline $[0{,}024;\ 0{,}06]$ &
    пуассонівські, інтервали $\exp(1/\lambda_i)$~хв \\ \hline
\end{kpitable}

Незалежні фактори -- лише $s$ і~$\gamma$, інтенсивності $\lambda_i=\gamma\lambda_i^0$ похідні;
окремо їх змінюють лише у верифікації (розділ~3) і пробній серії розділу~4.

\begin{kpitable}{|L{6.2cm}|c|L{7.2cm}|}{Параметри моделі\label{tab:params}}
  {\textbf{Параметр} & \textbf{Позначення} & \textbf{Значення}}
  Інтенсивності виконання A, B, C & $\mu_A, \mu_B, \mu_C$ &
    $0{,}25$; $0{,}2$; $0{,}1$~завд./хв, час $\exp(4)$, $\exp(5)$, $\exp(10)$~хв \\ \hline
  Обсяг оперативної памʼяті & $M$ & 2~умовні одиниці \\ \hline
  Потреба A, B, C у памʼяті & $v_A, v_B, v_C$ & 1; 1; 2~умовні одиниці \\ \hline
  Ємність черги & -- & необмежена \\ \hline
  Ймовірності вибору A і~B за рівних пріоритетів & -- & $0{,}5$; $0{,}5$ (лише стратегія~5) \\ \hline
  Час моделювання, тривалість перехідного періоду & $T_{\text{мод}}$, $T_{\text{п}}$ &
    обґрунтовуються в розділі~4 \\ \hline
\end{kpitable}

\subsection{Вихідні змінні моделі}

Вихідні змінні (таблиця~\ref{tab:outputs}; формули -- у розділі~2) визначають для кожного типу
і загалом (середні за всіма завданнями), бо підвищення пріоритету одного типу скорочує його
очікування за рахунок інших; «середній час обслуговування в системі» тлумачимо як час перебування.

\begin{kpitable}{|L{4.3cm}|c|c|L{6.9cm}|}{Вихідні змінні моделі\label{tab:outputs}}
  {\textbf{Змінна} & \textbf{Позначення} & \textbf{\makecell{Одиниці\\виміру}} & \textbf{Що спостерігається під час імітації}}
  Середній час перебування в системі & $W_i$, $W$ & хв &
    кількість завдань типу~$i$ у системі (середня за часом), надходжень \\ \hline
  Середня кількість завдань, що очікують обслуговування & $L_{q,i}$, $L_q$ & завд. &
    середня за часом довжина підчерги типу~$i$ \\ \hline
  Середній час очікування в черзі & $W_{q,i}$, $W_q$ & хв &
    довжина підчерги (середня за часом), кількість надходжень \\ \hline
  Коефіцієнт використання памʼяті (допоміжна) & $U$ & -- &
    середня за часом кількість зайнятих одиниць памʼяті${}/M$ \\ \hline
\end{kpitable}


\subsection{Висновки до розділу}

Умову конкретизовано припущеннями (пуассонівські потоки, не більше двох одночасних завдань,
резервування, без переривання), через перевантаження ($\rho=1{,}175$) інтенсивності варіюють
коефіцієнтом~$\gamma$; вхідні ($s$, $\gamma$) і вихідні ($W$, $L_q$, $W_q$) змінні відповідають меті.
